\documentclass[10pt,a4paper]{amsart}
\usepackage[margin=19mm]{geometry}
\usepackage{amsmath,amssymb,mathtools}
\usepackage[scr=rsfs,cal=euler]{mathalfa}
\usepackage{xfrac}
\usepackage[unicode,hidelinks,bookmarks=false]{hyperref}
\newcommand{\ie}{i.e.\ }
\newcommand{\cf}{cf.\ }
\newcommand{\resp}{respectively\ }
\newcommand{\Subsection}{\subsection}
\providecommand{\nobreakdash}{\nobreak\hbox{-}}
\setlength{\emergencystretch}{2em}
\multlinegap0pt
\usepackage{amssymb}
\newtheorem{lemma}{Lemma}
\title{Minimal simplicial spherical mappings with a given degree}
\author{Ksenia Apolonskaya, Oleg R. Musin}
\date{Source: arXiv:2511.10870v4 (CC BY 4.0)}
\begin{document}
\maketitle

\noindent\emph{Source and licence.} Minimal simplicial spherical mappings with a given degree. Version arXiv:2511.10870v4 (CC BY 4.0), distributed by arXiv under the Creative Commons Attribution 4.0 International licence, http://creativecommons.org/licenses/by/4.0/, as recorded in the arXiv licence metadata for this exact version and shown on the abstract page https://arxiv.org/abs/2511.10870v4. Full licence notice: \texttt{licenses/arxiv-2511.10870-CC-BY-4.0.txt}.\par\smallskip

\noindent\textit{Editorial scope.} Counted unit: the article's lemma lem:4 (source0.tex lines 207--210). The article's complete original proof of it is delivered with it (source0.tex lines 211--222, one \texttt{proof} environment of 12 lines). No mathematics is written, repaired or re-derived: the statement and the proof are the article's own lines, in the article's own notation, and no retained line is substituted or re-flowed.  No further source-local context is needed: every cross-reference of every retained line resolves inside the two delivered spans.  Nothing was omitted from any retained span, so the package carries no omission record.  No retained line contains a citation, so the wrapper carries no bibliography and no bibliography entry.  No retained line contains a figure environment, an image-inclusion command or drawing code, so no image is delivered and the package declares no asset dependency.  Editorial formatting only, in addition: an AMS \texttt{amsart} preamble that restates the article's own declarations and macros used by the retained lines, namely the environments \texttt{lemma} on separate counters, together with the standard packages those lines need.  The retained environments print in this document as Lemma 1 in order of appearance, whereas the article prints the same items with the numbering of its own full document.  The complete native source of the article as distributed in the arXiv e-print for this exact version is bundled unchanged as \texttt{sources/189072/source0.tex}.
\begin{lemma}
\label{lem:4}
$\lambda(3,4)=10$
\end{lemma}

\begin{proof}
We first prove that $\lambda(3,4) \geq 10$.\\
The simplex $\Delta^4$ has five vertices. Let $U_i$ denote the set of preimages of vertex $i$, for $i=1,\dots,5$. If $|U_i| \geq 2$ for all $i$, then $\lambda(3,4) \geq 10$.

Now suppose $|U_1| = 1$. Remove $U_1$ together with all vertices not connected to it and the corresponding edges. The result is a triangulation of $S^2$ with degree $4$, which requires at least $\lambda(2,4) = 10$ vertices. Adding back $U_1$ gives at least $11$ vertices.

We now construct an example with exactly $10$ vertices. Start with the standard simplex $u_1 u_2 u_3 u_4 u_5$. Consider the face $u_1 u_2 u_3 u_4$ and insert vertex $w_5$, replacing it with simplices $u_1 u_2 u_3 u_5, \dots, u_2 u_3 u_4 u_5$. Repeat this process for all other faces, adding vertices $w_1, w_2, w_3, w_4$.

The resulting triangulation has $10$ vertices and four preimages of $v_1 v_2 v_3 v_4$ (namely $w_1 v_2 v_3 v_4$, $v_1 w_2 v_3 v_4$, $v_1 v_2 w_3 v_4$, $v_1 v_2 v_3 w_4$), all negatively oriented. Permuting the vertex ordering in $v_1 v_2 v_3 v_4 v_5$ yields degree $d=4$.

Thus, $\lambda(3,4) \geq 10$ and $\lambda(3,4) \leq 10$, so $\lambda(3,4) = 10$.
\end{proof}

\end{document}
